\documentclass[oneside]{diretrizes} % Imprimi apenas frente da página.

\usepackage[alf, abnt-emphasize=bf, recuo=0cm, abnt-etal-cite=2, abnt-etal-list=0]{abntex2cite} % Criação de citações seguindo as normas ABNT. Os parâmetros especificam o estilo de citação a ser utilizado.
\usepackage[utf8]{inputenc} % Permite o uso de caracteres acentuados diretamente no código LaTeX.
\usepackage[T1]{fontenc} % Define a codificação da fonte para 8 bits, permitindo a correta exibição de caracteres especiais, como acentos.
\usepackage{graphicx} % Permite a inserção e manipulação de figuras no documento.
\usepackage{amsmath,amsthm,amsfonts,amssymb,dsfont,mathtools} % Pacotes para aprimorar a formatação e funcionalidade de fórmulas matemáticas, além de fornecer símbolos e fontes específicas para matemática.
\usepackage{booktabs} % Fornece comandos adicionais para criar tabelas de alta qualidade, com linhas horizontais mais agradáveis e espaçamento apropriado.
\usepackage{verbatim} % Permite que o texto seja interpretado exatamente como está escrito no documento, sem processamento especial. Útil para exibir código-fonte ou blocos de texto preformatados.
\usepackage{multirow, array} % Pacotes para trabalhar com tabelas que envolvem a combinação de várias linhas ou colunas.
\usepackage{indentfirst} % Indenta o primeiro parágrafo de cada seção. É uma prática comum em alguns estilos de formatação.
\usepackage{microtype} % Melhora a justificação do texto no documento, fazendo pequenos ajustes no espaçamento entre letras e palavras.
\usepackage[algoruled, portuguese]{algorithm2e} % Permite escrever algoritmos no documento com um estilo agradável, incluindo números de linha e estruturas de controle adequadas.
\usepackage{float} % Utilizado para criar ambientes "float" como figuras e tabelas, permitindo maior controle de posicionamento.
\usepackage{times} % Define a fonte do documento como Times.
\usepackage{afterpage} % Pacote que permite inserir uma página em branco após a página atual.
\usepackage{tocbibind} % Adiciona entradas para a bibliografia, índice e lista de figuras/tables ao índice do documento.
\usepackage{leading} % Controla o espaçamento entre linhas.
\usepackage{hyperref} % Fornece hiperlinks clicáveis no documento PDF gerado. Também é útil para criar referências internas e externas.
\newcolumntype{A}{>{\centering\arraybackslash}p{1.2cm}} % Define um novo tipo de coluna para tabelas com alinhamento centralizado.
\usepackage{url} % Permite a formatação e quebra de URLs longas dentro do texto.
\usepackage[bottom]{footmisc} % Coloca as notas de rodapé na parte inferior da página.
\usepackage{pdfpages} % Permite a inclusão de páginas de arquivos PDF externos no documento LaTeX.
\usepackage{pgffor} % Fornece comandos para realizar iterações, como loops, dentro do LaTeX. É útil para tarefas como repetir um bloco de código várias vezes.
\usepackage{enumitem} % Permite personalizar as listas (itemize, enumerate, description) no LaTeX. Ele oferece uma ampla variedade de opções de formatação e personalização para listas.

%%%%%%%%%%%%%%%%% Incluir o Sudoku %%%%%%%%%%%%%%%%%%
\usepackage{tikz} %utilizado para a criação de gráficos, diagramas e desenhos vetoriais
\newcounter{row}
\newcounter{col}

\newcommand{\setrow}[9]{
	\setcounter{col}{1}
	\foreach \n in {#1, #2, #3, #4, #5, #6, #7, #8, #9} {
		\edef\x{\value{col} - 0.5}
		\edef\y{9.5 - \value{row}}
		\node[anchor=center] at (\x, \y) {\n};
		\stepcounter{col}
	}
	\stepcounter{row}
}

\newcommand{\setrowfor}[9]{
	\setcounter{col}{1}
	\foreach \n in {#1, #2, #3, #4} {
		\edef\x{\value{col} - 0.5}
		\edef\y{4.5 - \value{row}}
		\node[anchor=center] at (\x, \y) {\n};
		\stepcounter{col}
	}
	\stepcounter{row}
}

%%%%%%%%%%%%%%%%%%%%%%%%%%%%%%%%%%%%%%%%%%%%%%%%%%%%%%%%%%%%%%para esconder a contagem de números basta colocar o * após o newtheorem, por exemplo \newtheorem*{teo}{Teorema}[] e ficará como Teorema.%

\newtheorem{teo}{Teorema}[]
\newtheorem{prop}{Proposição}[]
\newtheorem{lema}{Lema}[]
\newtheorem{coro}{Corolário}[]
\theoremstyle{definition}
\newtheorem{defi}{Definição}[]
\newtheorem{observacao}{Observação}[]
\newtheorem{ex}{Exemplo}[]
\newtheorem{sol}{Solução}[] 

% Inclui o preâmbulo do documento
\include{01-elementos-pre-textuais/preambulo}

% Define as cores dos links e informações do PDF
\makeatletter
\hypersetup{
    portuguese,
    colorlinks,
    linkcolor=black,
    citecolor=black,
    filecolor=black,
    urlcolor=black,
    breaklinks=true,
    pdftitle={\@title},
    pdfauthor={\@author},
    pdfsubject={\imprimirpreambulo},
    pdfkeywords={abnt, latex, abntex, abntex2}
}
\makeatother

% Redefinição de labels
\renewcommand{\algorithmautorefname}{Algoritmo}
\def\equationautorefname~#1\null{Equa\c c\~ao~(#1)\null}

% Cria o índice remissivo
\makeindex

% Início do documento
\begin{document}

    % Retira espaço extra obsoleto entre as frases.
    \frenchspacing

    % Elementos pré textuais
    \pretextual
    \include{./01-elementos-pre-textuais/capa}              % Capa
    \include{./01-elementos-pre-textuais/folhaRosto}        % Folha de rosto
    \begin{center}
    
    \end{center}
    \include{./01-elementos-pre-textuais/folhaAprovacao}    % Folha de aprovação
    \include{./01-elementos-pre-textuais/dedicatoria}       % Dedicatória
    \include{./01-elementos-pre-textuais/agradecimentos}    % Agradecimentos
    \include{./01-elementos-pre-textuais/epigrafe}          % Epígrafe
    \include{./01-elementos-pre-textuais/resumoPt}          % Resumo na língua vernácula
    \include{./01-elementos-pre-textuais/resumoEn}          % Resumo em língua estrangeira
    \include{./01-elementos-pre-textuais/listaFiguras}      % Lista de figuras
    \include{./01-elementos-pre-textuais/listaTabelas}      % Lista de tabelas
   
    \sumario

    % Elementos textuais
    \textual
    \include{./02-elementos-textuais/introducao}            % Introdução
    \include{./02-elementos-textuais/Cap1}           % Disposições
    \include{./02-elementos-textuais/Cap2}  			% Estrutura    
    \include{./02-elementos-textuais/Cap3}           	% Ilustrações
    \include{./02-elementos-textuais/Conclusões}
    
    % Elementos pós textuais
    \postextual
    % Referências
    \include{./03-elementos-pos-textuais/Referências}
    
    %%%%%%%%%%%%%%%%%%%%%%%%%%%%%%%%%%%%%%%%%%%%%%%%%%%%%%%
    % IMPORTANTE!!! Cada referência deve ter seu nome posto no modelo ``\nocite{nomedareferencia}'', como abaixo.
    % Além disso, DEVE-SE LEMBRAR DE CITAR CADA REFERÊNCIA NO CORPO DO TEXTO. Por definição, ``Referências''
    % é uma lista de bibliografias onde se encontram afirmações e indicações de conteúdo do texto, e serve para o leitor
    % consultar e conferir a existência/veracidade destas afirmações.
    %%%%%%%%%%%%%%%%%%%%%%%%%%%%%%%%%%%%%%%%%%%%%%%%%%%%%%%
    \nocite{autor}
    \nocite{referência}
    \nocite{internet}
    \nocite{net}
    %%%%%%%%%%%%%%%%%%%%%%%%%%%%%%%%%%%%%%%%%%%%%%%%%%%%%%%

\end{document}
